\documentclass[12pt]{article}
\usepackage[a4paper,margin=32mm]{geometry}
\usepackage[T1]{fontenc}
\usepackage{lmodern}
\usepackage{amsmath,amssymb,amsthm,mathtools}
\usepackage{booktabs,needspace}
\usepackage[final]{microtype}
\usepackage[hidelinks]{hyperref}
\hypersetup{
  pdftitle={Affine Dual Jacobi--Trudi Identities for Classical Hall--Littlewood Polynomials via Complementary Minors},
  pdfauthor={Helen W. J. Zhang and Chengdong Zhao},
  pdfsubject={Affine dual Jacobi--Trudi identities},
  pdfkeywords={Hall--Littlewood polynomials, complementary minors, Jacobi--Trudi identities}
}
\newtheorem{theorem}{Theorem}[section]
\newtheorem{maintheorem}{Theorem}
\renewcommand{\themaintheorem}{\Alph{maintheorem}}
\newtheorem{proposition}[theorem]{Proposition}
\newtheorem{lemma}[theorem]{Lemma}
\newtheorem{corollary}[theorem]{Corollary}
\numberwithin{equation}{section}
\newcommand{\Q}{\mathbb Q}
\newcommand{\Z}{\mathbb Z}
\newcommand{\dote}{\dot e}
\newcommand{\ddote}{\ddot e}
\newcommand{\dotE}{\dot E}
\newcommand{\ddotE}{\ddot E}
\newcommand{\ph}{\phi}
\newcommand{\ps}{\psi}
\newcommand{\Jaff}{\mathfrak J}
\newcommand{\diag}{\operatorname{diag}}
\newcommand{\latticeeven}{\Z_{\mathrm{even}}^k}
\newcommand{\eb}{\mathbf e}
\setlength{\emergencystretch}{2em}
\setlength{\jot}{4pt}
\clubpenalty=10000
\widowpenalty=10000
\displaywidowpenalty=10000
\raggedbottom
% Title
\begin{document}
\begin{center}
\LARGE
Affine Dual Jacobi--Trudi Identities\\[3pt]
for Classical Hall--Littlewood Polynomials\\[3pt]
via Complementary Minors
\end{center}
% GITHUB-EDITION:AUTHORS-BEGIN
\begin{center}
\large Helen W.~J. Zhang$^{1}$ \qquad Chengdong Zhao$^{2}$\\[6pt]
\small
$^{1}$School of Mathematics, Hunan University\\
\href{mailto:helenzhang@hnu.edu.cn}{\texttt{helenzhang@hnu.edu.cn}}\\[5pt]
$^{2}$School of Mathematics and Statistics, Central South University\\
\href{mailto:cdzhao@csu.edu.cn}{\texttt{cdzhao@csu.edu.cn}}
\end{center}
% GITHUB-EDITION:AUTHORS-END
\medskip

\begin{abstract}
We prove the six affine dual Jacobi--Trudi identities conjectured by
Warnaar for rectangular Hall--Littlewood polynomials of types $B_n$,
$C_n$, and $BC_n$, together with his additional type-$B_n$ identity at
$s=t$. We give a uniform proof by adapting the complementary-minor
method used by Aitken for symmetric-function duality.
\end{abstract}

\medskip
\noindent\textbf{Keywords.} Hall--Littlewood polynomials; Jacobi--Trudi
identities; Weyl characters; complementary minors; affine root systems.

\smallskip
\noindent\textbf{2020 Mathematics Subject Classification.}
05E05; 05E10; 15A15; 17B67.

\section{Introduction and main identities}

Classical dual Jacobi--Trudi identities express rectangular Weyl
characters of rank $n$ as determinants of order $k$, where $k$ is the
width of the rectangle.  Warnaar proposed affine analogues for classical
Hall--Littlewood polynomials in \cite[Conjectures~1.2--1.4]{Warnaar2026},
and a further type-$B_n$ identity in \cite[Section~6]{Warnaar2026}.
We prove these identities, with one necessary correction at width zero.
Warnaar established the $k=1$ cases of the six numbered conjectures and
the $t=1$ specializations (1.6a) and (1.8)
\cite[Lemmas~4.1--4.4]{Warnaar2026}; his type-$A$ analogue is
\cite[Theorem~1.1]{Warnaar2026}.  The applications to affine and
$q,t$-Rogers--Ramanujan identities developed there use the classical-type
formulas proved here.

For $x=(x_1,\ldots,x_n)$, define
\begin{equation}\label{eq:elementary}
\begin{aligned}
 \dotE(u)&=\prod_{a=1}^n(1+ux_a)(1+ux_a^{-1})
          =\sum_{r=0}^{2n}\dote_r u^r,\\
 \ddotE(u)&=(1+u)\dotE(u)=\sum_{r=0}^{2n+1}\ddote_r u^r.
\end{aligned}
\end{equation}
Both coefficient families are extended by zero to all integer indices.
Thus
\[
 \dote_{2n-r}=\dote_r,\qquad
 \ddote_{2n+1-r}=\ddote_r,\qquad
 \ddote_r=\dote_r+\dote_{r-1}.
\]
For $1\le i\le k$ and $m\in\Z$, use the row kernels:
\begin{align*}
 \Lambda_i(m)&=\dote_{n-i+m}-\dote_{n+i+m},&
 h_i(m)&=\dote_{n-i+m}+\dote_{n+i+m-1},\\
 h_i^+(m)&=\dote_{n-i+m}+\dote_{n+i+m-2},&
 \widehat h_i(m)&=\ddote_{n-i+m}+\ddote_{n+i+m-1}.
\end{align*}
For weights $w_j(y)$, kernels $F_i(m)$, and
$\mathcal L\subseteq\Z^k$, put
\begin{equation}\label{eq:J}
 \Jaff_K[w,F;\mathcal L]
 :=\sum_{(y_1,\ldots,y_k)\in\mathcal L}
       \det_{1\le i,j\le k}[w_j(y_i)F_i(j-Ky_i)].
\end{equation}
In this notation an expression in $j,y$ in the weight slot specifies
$w_j(y)$.  All sums below are finite by zero extension.  Write
$\latticeeven=\{y\in\Z^k:y_1+\cdots+y_k\in2\Z\}$.
The empty determinant and the sum indexed by the empty vector both have
value $1$.  When $k=0$, no entry or weight in \eqref{eq:J} is evaluated.
The coefficient field is $\Q(t,x_1,\ldots,x_n)$, except in the
two square-root cases, where it is $\Q(\tau,x_1,\ldots,x_n)$ and
$t^{r/2}:=\tau^r$ for $r\in\Z$.

\Needspace{17\baselineskip}
\begin{maintheorem}\label{thm:A}
For $n\ge1$, the following identities hold for $k\ge0$, except that
\eqref{eq:main-Bplus} requires $k\ge1$:
\begin{align}
 P_{(k^n)}^{B_n}(x;t,0)
 &=\Jaff_{2k+1}[(-1)^y t^{(k+\frac12)y^2-(j-\frac12)y},h;\Z^k],
 \label{eq:main-B0}\\
 P_{(k^n)}^{B_n}(x;\tau^2,-\tau)
 &=\Jaff_{2k}[\tau^{2ky^2-(2j-1)y},h;\Z^k],
 \label{eq:main-Bminus}\\
 P_{(k^n)}^{B_n}(x;t,1)
 &=\tfrac12\Jaff_{2k}[(-1)^y t^{ky^2-(j-1)y},h^+;\Z^k],
 \label{eq:main-Bplus}\\
 P_{(k^n)}^{C_n}(x;t,0)
 &=\Jaff_{2k+2}[t^{(k+1)y^2-jy},\Lambda;\Z^k],
 \label{eq:main-C0}\\
 P_{(k^n)}^{C_n}(x;t,t)
 &=\Jaff_{2k}[t^{ky^2-jy},\Lambda;\latticeeven],
 \label{eq:main-Ct}\\
 P_{(k^n)}^{BC_n}(x;\tau^2,-\tau,0)
 &=\Jaff_{2k+1}[(-1)^y\tau^{(2k+1)y^2-2jy},\Lambda;\Z^k].
 \label{eq:main-BC}
\end{align}
At width zero the normalized value in \eqref{eq:main-Bplus} is
$P_{\varnothing}^{B_n}(x;t,1)=1$.
\end{maintheorem}

These are, in order, Conjectures~1.2(a)--(c), 1.3(a)--(b), and 1.4 of
\cite{Warnaar2026}.  The literal $k=0$ assertion of Conjecture~1.2(c)
would give $1/2$ on the right of \eqref{eq:main-Bplus}; its correct
normalized endpoint is therefore separate from the positive-width
identity.

\Needspace{8\baselineskip}
\begin{maintheorem}\label{thm:B}
For $n,k\ge1$,
\begin{equation}\label{eq:main-Bt}
 P_{(k^n)}^{B_n}(x;t,t)
 =\tfrac12\Jaff_{2k}[(-1)^y t^{ky^2-(j-1)y},\widehat h;\Z^k].
\end{equation}
The normalized width-zero value is $P_{\varnothing}^{B_n}(x;t,t)=1$.
\end{maintheorem}

Theorem~\ref{thm:B} is the additional, unnumbered identity in
\cite[Section~6]{Warnaar2026}.  Our Hall--Littlewood conventions are
those of \cite[(2.16)--(2.17)]{Warnaar2026}; for the type-$BC$ Weyl
sum and polynomiality, see also
\cite[Definition~2.4 and Theorems~2.5--2.6]{Venkateswaran2013}.

The proof uses complementary minors, a familiar source of dual
symmetric-function determinants \cite{Aitken1931,McDowell2023}.
The corresponding proportionality of complementary maximal minors is
also the basic Gale-dual relation; see
\cite[Lemma~2.10]{MullerEtAl2016}.  We apply it to the kernel pairs
$\Phi D=0$ and $V^{\mathsf T}C=0$.
Here $\det(\Phi V)/\Delta$ is a finite Weyl representation, while
$D^{\mathsf T}C$ is the affine matrix.  The coefficient kernels $D$
come from \eqref{eq:elementary}; the affine kernels $C$ come from
reflections of weighted residue classes.  Explicit complementary blocks
fix all signs and constants.

Section~\ref{sec:finite} derives the finite Weyl determinants.
Section~\ref{sec:transfer} gives the normalized transfer and the three
coefficient kernels.  Section~\ref{sec:ordinary} treats four ordinary
cases in two pairs.  Section~\ref{sec:exceptional} treats the fixed
residue, fixed center, and adjacent residues responsible for the three
half-determinant formulas.  Width zero and the classical specializations
are collected in Section~\ref{sec:boundary}.


\Needspace{10\baselineskip}
\section{Finite Weyl determinants}\label{sec:finite}

We abbreviate the seven polynomials by $P^\star_{n,k}$, where
$\star=C_0,C_t,BC,B_0,B_-,B_+,B_t$ denotes the type and boundary
parameter in Theorems~\ref{thm:A} and~\ref{thm:B}.
For $(a;t)_m=\prod_{r=0}^{m-1}(1-at^r)$, set
\[
 V_n(t)=\prod_{r=1}^n(1+t+\cdots+t^{r-1}).
\]
The rectangular Weyl-sum definition, including width zero, is
\begin{equation}\label{eq:weyl}
 P^\star_{n,k}
 =\frac{1}{V_n(t)\nu_\star(k)}
 \sum_{w\in W_n}w\left[
 \prod_{a=1}^n x_a^{-k}\rho_\star(x_a)
 \prod_{a<b}\frac{(tx_a-x_b)(1-tx_ax_b)}
                    {(x_a-x_b)(1-x_ax_b)}\right],
\end{equation}
where $W_n=S_n\ltimes\{\pm1\}^n$ acts by permutations and inversions,
$\nu_\star(k)=1$ for $k\ge1$, and the remaining data are as follows.
In the $BC$ and $B_-$ rows, $t=\tau^2$.
\begin{center}
\renewcommand{\arraystretch}{1.35}
\begin{tabular}{@{}lll@{}}
\toprule
Case $\star$ & $\rho_\star(z)$ & $\nu_\star(0)$\\
\midrule
$C_0$ & $(1-z^2)^{-1}$ & $1$\\
$C_t$ & $(1-tz^2)/(1-z^2)$ & $(-t;t)_n$\\
$BC$ & $(1+\tau z)/(1-z^2)$ & $1$\\
$B_0$ & $(1-z)^{-1}$ & $1$\\
$B_-$ & $(1+\tau z)/(1-z)$ & $(\tau;\tau^2)_n$\\
$B_+$ & $1$ & $(-1;t)_n$\\
$B_t$ & $(1-tz)/(1-z)$ & $(-t;t)_n$\\
\bottomrule
\end{tabular}
\end{center}
Indeed, the symmetric-group normalization of either $(k^n)$, $k>0$,
or $\varnothing$ is $V_n(t)$; only the boundary Pochhammer factor
changes at width zero.

The standard symmetrizer identity
\[
 \sum_{\sigma\in S_n}
 \prod_{a<b}\frac{tz_{\sigma(a)}-z_{\sigma(b)}}
                  {z_{\sigma(a)}-z_{\sigma(b)}}=V_n(t)
\]
follows by alternation and comparison of degrees; setting $z_n=0$
gives the recurrence $V_n=(1+t+\cdots+t^{n-1})V_{n-1}$.
Since the other factors in \eqref{eq:weyl} are symmetric, it reduces
that expression to
\begin{equation}\label{eq:signsum}
 P^\star_{n,k}
 =\frac{1}{\nu_\star(k)}
   \sum_{\eta\in\{\pm1\}^n}
   \prod_a u_a^{-k}\rho_\star(u_a)
   \prod_{a<b}\frac{1-tu_au_b}{1-u_au_b},
 \qquad u_a=x_a^{\eta_a}.
\end{equation}

\subsection{Alternants and finite foldings}

Write $N_t(z)=\prod_{a<b}(z_a-z_b)(1-tz_az_b)$ and put
\[
 f_1(z)=z^{n-1},\qquad
 f_r(z)=z^{n-r}+t^{r-1}z^{n+r-2}\quad(2\le r\le n).
\]
The finite numerator calculations reduce to three alternants.

\begin{lemma}
With rows indexed by $1\le a\le n$ and columns by $1\le r\le n$,
\begin{align}
 \det[f_r(z_a)]&=N_t(z),\label{eq:alternant-basic}\\
 \det[z_a^{n-r}-t^rz_a^{n+r}]
 &=\prod_a(1-tz_a^2)N_t(z),\label{eq:alternant-Ct}\\
 \det[z_a^{n-r}+\tau^{2r-1}z_a^{n+r-1}]
 &=\prod_a(1+\tau z_a)N_{\tau^2}(z).
 \label{eq:alternant-Bminus}
\end{align}
\end{lemma}

\begin{proof}
Each determinant is alternating.  On setting $z_b=(tz_a)^{-1}$,
rows $a,b$ are proportional, with multipliers
$t^{1-n}z_a^{2-2n}$, $-t^{-n}z_a^{-2n}$, and
$\tau^{1-2n}z_a^{1-2n}$, respectively; in the last expression
$t=\tau^2$.  This gives the factors $1-tz_az_b$.
The second determinant has row factor $1-tz_a^2$, and the third
vanishes at $z_a=-\tau^{-1}$.  Divisibility can be checked over the
fraction field in the remaining variables and returned to the polynomial
ring by Gauss's lemma.  The factors are pairwise coprime and,
together with $z_a-z_b$, have
total degrees
\[
 3\binom n2,\qquad 2n+3\binom n2,\qquad n+3\binom n2,
\]
which are the respective upper bounds on the determinant degrees.
In each case the lowest-degree homogeneous part is
$\det[z_a^{n-r}]=\prod_{a<b}(z_a-z_b)$, fixing the constant at $1$.
\end{proof}

Set $\ph_m(z)=z^m-z^{-m}$, $\ps_m(z)=z^m-z^{1-m}$, and
$C_m(z)=z^m+z^{-m}$.
For $L=n+k$ and standard basis $\eb_1,\ldots,\eb_L$, use the foldings
\[
\begin{aligned}
 [m]_C&=\begin{cases}\eb_m&m>0,\\0&m=0,\\-\eb_{-m}&m<0,\end{cases}
 & [m]_B&=\begin{cases}\eb_m&m\ge1,\\-\eb_{1-m}&m\le0,\end{cases}\\
 [m]_+&=\begin{cases}\eb_m&m\ge2,\\2\eb_1&m=1,\\\eb_{2-m}&m\le0.\end{cases}
\end{aligned}
\]
Only indices whose folded rows lie in $\{1,\ldots,L\}$ will occur.
The evaluation matrices are
\[
\begin{gathered}
 \Phi^C_{a,p}=\ph_p(x_a),\qquad \Phi^B_{a,p}=\ps_p(x_a),\\[2pt]
 \Phi^+_{a,p}=\begin{cases}
 1,&p=1,\\
 C_{p-1}(x_a),&p\ge2.
 \end{cases}
\end{gathered}
\]
Thus $\Phi^C[m]_C=\ph_m(x)$, $\Phi^B[m]_B=\ps_m(x)$,
and $\Phi^+[m]_+=C_{m-1}(x)$, with vector evaluation understood.
In particular, $[1]_+=2\eb_1$ is forced by $C_0=2$.

The three initial minors are the Weyl denominators
\begin{align}
 \Delta_C(z)&=\det[\ph_r(z_a)]
 =(-1)^n\prod_a z_a^{-n}(1-z_a^2)N_1(z),\notag\\
 \Delta_B(z)&=\det[\ps_r(z_a)]
 =(-1)^n\prod_a z_a^{-(n-1)}(1-z_a)N_1(z),\label{eq:denominators}\\
 \Delta_+(z)&=\det[\chi_r(z_a)]
 =\prod_a z_a^{-(n-1)}N_1(z),\notag
\end{align}
where $\chi_1=1$ and $\chi_r=C_{r-1}$ for $r\ge2$.
These evaluations are immediate Vandermonde determinants in $z+z^{-1}$:
$\ph_r/(z-z^{-1})$, $\ps_r/(z-1)$, and $\chi_r$ are monic
polynomials of degree $r-1$.  Inversion of a row multiplies the three
denominators by $-1$, $-z^{-1}$, and $1$, respectively.

\subsection{The finite column matrices}

For $\diamond\in\{C,B,+\}$, define
\[
\begin{aligned}
 u_1^\diamond&=[k+1]_\diamond,\qquad
 \widetilde u_1^\diamond=[k]_\diamond,\\
 u_r^\diamond&=[k+r]_\diamond+t^{r-1}[k+2-r]_\diamond
 \quad(r\ge2),\\
 \widetilde u_r^\diamond&=[k+r-1]_\diamond+t^{r-1}[k+1-r]_\diamond
 \quad(r\ge2).
\end{aligned}
\]
The notation $\widetilde u_r^\diamond$ is used only for $\diamond=C,B$.
Let $V^\star=(v_1^\star,\ldots,v_n^\star)$ be given by
Table~\ref{tab:finite-columns}; every folded row there is at most $k+r$
in column $r$.

\begin{table}[htbp]
\centering
\renewcommand{\arraystretch}{1.35}
\begin{tabular}{@{}llll@{}}
\toprule
$\star$ & $\Phi$ & $v_r^\star$ & $\Delta$\\
\midrule
$C_0$ & $\Phi^C$ & $u_r^C$ & $\Delta_C$\\
$C_t$ & $\Phi^C$ & $[k+r]_C-t^r[k-r]_C$ & $\Delta_C$\\
$BC$ & $\Phi^C$ & $u_r^C+\tau\widetilde u_r^C$ & $\Delta_C$\\
$B_0$ & $\Phi^B$ & $u_r^B$ & $\Delta_B$\\
$B_-$ & $\Phi^B$ & $[k+r]_B+\tau^{2r-1}[k+1-r]_B$ & $\Delta_B$\\
$B_+$ & $\Phi^+$ & $u_r^+$ & $\Delta_+$\\
$B_t$ & $\Phi^B$ & $u_r^B-t\widetilde u_r^B$ & $\Delta_B$\\
\bottomrule
\end{tabular}
\caption{Finite column data.  In the $BC$ and $B_-$ rows, $t=\tau^2$.}
\label{tab:finite-columns}
\end{table}

\begin{proposition}\label{prop:finite}
For $n,k\ge1$, the data of Table~\ref{tab:finite-columns} satisfy
\[
 P^\star_{n,k}=\frac{\det(\Phi V^\star)}{\Delta}.
\]
Moreover, $\det V^\star_H=1$ for $H=\{k+1,\ldots,k+n\}$.
\end{proposition}

\begin{proof}
Multiply \eqref{eq:signsum} by the relevant denominator in
\eqref{eq:denominators}.  In the $C_0$ case,
\eqref{eq:alternant-basic} gives row functions
$z^{-n-k}f_r(z)$.  Summing over inversions folds each row by
$g(z)\mapsto g(z)-g(z^{-1})$, and these functions give
$-\Phi^C u_r^C$.  The $n$ column signs cancel the factor $(-1)^n$
in $\Delta_C$.  The $BC$ calculation multiplies the same row
functions by $1+\tau z$, giving
$-\Phi^C(u_r^C+\tau\widetilde u_r^C)$.
For $C_t$, replace $f_r$ by $z^{n-r}-t^rz^{n+r}$ using
\eqref{eq:alternant-Ct}; the folded column is
$-\Phi^C([k+r]_C-t^r[k-r]_C)$.

For the shifted folding use $z^{-(n+k-1)}f_r(z)$ and
$g(z)\mapsto g(z)-zg(z^{-1})$.  This gives
$-\Phi^B u_r^B$ in case $B_0$; multiplying by $1-tz$ gives
$-\Phi^B(u_r^B-t\widetilde u_r^B)$ in case $B_t$.
Equation~\eqref{eq:alternant-Bminus} gives the $B_-$ column in
exactly the same way.  Again the column signs cancel $(-1)^n$.

In case $B_+$, $\Delta_+$ is invariant under each inversion.
Use \eqref{eq:alternant-basic} and the row functions
$z^{-(n+k-1)}f_r(z)$, now with
$g(z)\mapsto g(z)+g(z^{-1})$.  The first folded column is $C_k$,
and the others are
$C_{k+r-1}+t^{r-1}C_{k+1-r}$, exactly $\Phi^+u_r^+$.
This also checks the normalization of the first column without
introducing an additional even-orthogonal alternant.

Finally, in every case the leading bracket places a $1$ in row $k+r$
of column $r$.  For $k\ge1$ all other folded rows are strictly
smaller, including those arising from nonpositive indices.  Hence
$V^\star_H$ is unit upper triangular.
\end{proof}


\section{Complementary kernels and transfer}\label{sec:transfer}

Throughout Sections~\ref{sec:transfer}--\ref{sec:exceptional}, assume
$n,k\ge1$ and write
\[
\begin{gathered}
 L=n+k,\qquad Q_0=\{1,\ldots,n\},\qquad J=\{n+1,\ldots,n+k\},\\
 H=\{k+1,\ldots,k+n\},\qquad I=\{1,\ldots,k\},\\
 \omega=(-1)^{k(k+1)/2},\qquad \zeta=(-1)^{nk}\omega.
\end{gathered}
\]
Subsets are always ordered increasingly.  A subscript on an
$n\times L$ matrix selects columns; on a matrix with $L$ rows it
selects rows.

\subsection{The normalized transfer}

\begin{theorem}
Let $F$ be a field, and let $X\in F^{n\times L}$,
$Y\in F^{L\times n}$, and $U,Z\in F^{L\times k}$ satisfy $XU=0$ and $Y^{\mathsf T}Z=0$.
If $\det X_{Q_0}$, $\det U_J$, $\det Y_H$, and $\det Z_I$ are
nonzero, then
\[
 \det(XY)=
 \frac{(-1)^{nk}\det X_{Q_0}\det Y_H}
      {\det U_J\det Z_I}\det(U^{\mathsf T}Z).
\]
\end{theorem}

\begin{proof}
The nonzero minors give maximal rank, hence
$\operatorname{col}(U)=\ker X$.  The forms
\[
\begin{aligned}
 (v_1,\ldots,v_n)&\longmapsto\det[Xv_1\ \cdots\ Xv_n],\\
 (v_1,\ldots,v_n)&\longmapsto\det[U\ v_1\ \cdots\ v_n]
\end{aligned}
\]
induce nonzero top forms on $F^L/\operatorname{col}(U)$ and are
therefore proportional.  For $Q=\{q_1<\cdots<q_n\}$ set
$\epsilon(Q)=(-1)^{\sum_r(q_r-r)}$.  Expansion along the last $n$
columns gives
\[
 \det[U\ \eb_{q_1}\ \cdots\ \eb_{q_n}]
 =(-1)^{nk}\epsilon(Q)\det U_{Q^c}.
\]
Consequently, for constants $\alpha,\beta$ independent of $Q$,
\[
 \det X_Q=\alpha\epsilon(Q)\det U_{Q^c},\qquad
 \det Y_Q=\beta\epsilon(Q)\det Z_{Q^c},
\]
where the second relation follows from $Y^{\mathsf T}Z=0$.
Cauchy--Binet gives
\[
 \det(XY)=\alpha\beta
 \sum_{|Q|=n}\det U_{Q^c}\det Z_{Q^c}
 =\alpha\beta\det(U^{\mathsf T}Z).
\]
Evaluating the first relation at $Q_0$ and the second at $H$, for
which $\epsilon(Q_0)=1$ and $\epsilon(H)=(-1)^{nk}$, gives the result.
\end{proof}

In the finite representations, $\det\Phi_{Q_0}=\Delta$ and
$\det V_H=1$.  Hence matrices $D,C$ satisfying
$\Phi D=0$, $V^{\mathsf T}C=0$, and having nonzero minors $D_J,C_I$
give the single normalization formula
\begin{equation}\label{eq:normalized-transfer}
 P^\star_{n,k}
 =\frac{(-1)^{nk}}{\det D_J\det C_I}
   \det(D^{\mathsf T}C).
\end{equation}
Only the two indicated minors are needed in each application.

\subsection{Three coefficient kernels}

Palindromicity gives
\begin{equation}\label{eq:coefficient-reflections}
 \Lambda_i(-m)=-\Lambda_i(m),\qquad
 h_i(1-m)=h_i(m),\qquad h_i^+(2-m)=h_i^+(m).
\end{equation}
In particular,
\[
 h_i(m)=\dote_{n-i+m}+\dote_{n-i-m+1},\qquad
 h_i^+(m)=\dote_{n-i+m}+\dote_{n-i-m+2}.
\]
Define $L\times k$ matrices by
\begin{equation}\label{eq:D}
 D^C_{p,i}=(-1)^p\Lambda_i(p),\qquad
 D^B_{p,i}=(-1)^ph_i(p),\qquad
 D^+_{p,i}=(-1)^ph_i^+(p).
\end{equation}

\begin{lemma}\label{lem:coefficient-kernels}
The three pairs satisfy $\Phi^C D^C=0$, $\Phi^B D^B=0$, and
$\Phi^+D^+=0$.  Their complementary minors are
$\det D^C_J=\det D^B_J=\zeta$ and $\det D^+_J=2\zeta$.
\end{lemma}

\begin{proof}
The generating identities are
\begin{align*}
 \sum_m\Lambda_i(m)z^m&=z^{-n}\dotE(z)(z^i-z^{-i}),\\
 \sum_m h_i(m)z^m&=z^{-n}\dotE(z)(z^i+z^{1-i}),\\
 \sum_m h_i^+(m)z^m&=z^{-n}\dotE(z)(z^i+z^{2-i}).
\end{align*}
They vanish at $z=-x_a$.  Pair the integer indices by
$m\leftrightarrow-m$, $m\leftrightarrow1-m$, and
$m\leftrightarrow2-m$, respectively.  The first two paired sums are
$\sum_{p\ge1}(-1)^p\Lambda_i(p)\ph_p(x_a)$ and
$\sum_{p\ge1}(-1)^ph_i(p)\ps_p(x_a)$.
In the third, the fixed index $m=1$ gives, after division by $x_a$,
\[
 -h_i^+(1)+\sum_{p\ge2}(-1)^ph_i^+(p)C_{p-1}(x_a)=0.
\]
All three kernels vanish for $p>L$, giving the asserted kernel
relations.  In particular, the first row of $D^+$ has no extra factor $2$.

For $1\le r,i\le k$, the last $k$ rows satisfy
\begin{align*}
 D^C_{n+r,i}=D^B_{n+r,i}&=(-1)^{n+r}\dote_{i-r},\\
 D^+_{n+r,i}&=(-1)^{n+r}
       (\dote_{i-r}+\dote_{2-i-r}).
\end{align*}
These blocks are upper triangular.  The first two have diagonal
$(-1)^{n+r}$; the third has the same diagonal except that its first
entry is $2(-1)^{n+1}$, giving the stated minors.
\end{proof}

\subsection{Folding and finite unfolding}

We will fold scalar sequences $b_j(m)$, $m\in\Z$, into affine
matrices.  The three conventions are
\begin{align}
 \mathcal C^C(b)_{p,j}&=b_j(p)-b_j(-p),\notag\\
 \mathcal C^B(b)_{p,j}&=(-1)^p\bigl(b_j(p)+b_j(1-p)\bigr),
 \label{eq:affine-foldings}\\
 \mathcal C^+(b)_{p,j}&=
 \begin{cases}
 -b_j(1),&p=1,\\
 (-1)^p\bigl(b_j(p)+b_j(2-p)\bigr),&2\le p\le L.
 \end{cases}\notag
\end{align}
The sequences need not have finite support; the products below do.

\begin{lemma}
For the matrices in \eqref{eq:D} and \eqref{eq:affine-foldings},
\begin{align}
 ((D^C)^{\mathsf T}\mathcal C^C(b))_{ij}
 &=\sum_{m\in\Z}(-1)^m\Lambda_i(m)b_j(m),\label{eq:pairing-C}\\
 ((D^B)^{\mathsf T}\mathcal C^B(b))_{ij}
 &=\sum_{m\in\Z}h_i(m)b_j(m),\label{eq:pairing-B}\\
 ((D^+)^{\mathsf T}\mathcal C^+(b))_{ij}
 &=\sum_{m\in\Z}h_i^+(m)b_j(m).\label{eq:pairing-plus}
\end{align}
\end{lemma}

\begin{proof}
For the first identity, pair $m=p$ with $m=-p$, using
\eqref{eq:coefficient-reflections} and $\Lambda_i(0)=0$.
For the second, pair $p\ge1$ with $1-p\le0$.
For the third, the row $p=1$ contributes $h_i^+(1)b_j(1)$, while
$p\ge2$ pairs with $2-p\le0$.  Thus the fixed center is counted
once, not twice.  The vanishing for $p>L$ established above justifies
each unfolding.
\end{proof}

\Needspace{10\baselineskip}
A common bound controls all subsequent row sums.  If $K>0$ and
$\mathcal E_r=0$ outside $[0,N]$, a row with entries
\[
 \mathcal E_{n-i+j-Ky}
       +\sigma\mathcal E_{n+i+j-Ky-\delta},
 \qquad 1\le j\le k,\quad \delta\in\{0,1,2\},
\]
vanishes unless
\[
 \left\lceil\frac{n-i+1-N}{K}\right\rceil
 \le y\le
 \left\lfloor\frac{n+i+k-\delta}{K}\right\rfloor.
\]
Indeed, outside this interval every index is above $N$ or below $0$;
here $2i-\delta\ge0$.  Take $N=2n$, or $2n+1$ for
$\widehat h$.  Thus, if $A_{ij}=\sum_y w_j(y)F_i(j-Ky)$, finite row
multilinearity gives $\det A=\Jaff_K[w,F;\Z^k]$.
The even-sublattice case will use a parity projection instead.

Every weight has exponent
$\frac K2y^2-(j-\eta)y$, with
$\eta\in\{0,\tfrac12,1\}$,
$K/2\ge k-\eta$, and $j-\eta\ge0$.
It is nonnegative for every integer $y$.  In case $B_0$ it is an
integer because $y^2-y$ is even; in the square-root cases twice this
exponent is the exponent of $\tau$.  The resulting affine expressions
are therefore polynomial in $t$ or $\tau$, even though reflection
relations will be proved in the corresponding rational function fields.


\section{The four ordinary affine kernels}\label{sec:ordinary}

Four specializations have no doubled complementary block.  The two
unshifted foldings share one reflection, and the two shifted foldings
share another.

\subsection{The pair \texorpdfstring{$C_0$ and $BC$}{C0 and BC}}

\begin{proposition}
Equations~\eqref{eq:main-C0} and~\eqref{eq:main-BC} hold for $n,k\ge1$.
\end{proposition}

\begin{proof}
Take $K=2k+2$ for $C_0$ and $K=2k+1$ for $BC$, with $t=\tau^2$
in the latter case.  For $1\le j\le k$ define
\[
\begin{aligned}
 b_j(m)&=
 \begin{cases}
 t^{Kq^2/2+jq},&m=j+Kq,\quad q\in\Z,\\
 0,&m\not\equiv j\pmod K,
 \end{cases}\\
 c_j(m)&=b_j(m)-b_j(-m),
\end{aligned}
\]
and let $C=\mathcal C^C(b)$.  Shifting $q$ gives
$b_j(m+K)=t^{m+K/2}b_j(m)$, hence
\begin{equation}\label{eq:ordinary-C-reflection}
 c_j(-m)=-c_j(m),\qquad
 c_j(m)+t^{m-K/2}c_j(K-m)=0.
\end{equation}
The second relation follows by applying the quasiperiodicity to
$b_j(K-m)$ and $b_j(m-K)$; it holds also outside the residue classes.

For $C_0$, the relation at $m=k+1$ gives $c_j(k+1)=0$, and at
$m=k+r$ it gives
$c_j(k+r)+t^{r-1}c_j(k+2-r)=0$.  Since the odd extension agrees with
the bracket $[m]_C$, these are precisely
$(V^{C_0})^{\mathsf T}C=0$.
For $BC$, the first column gives
$c_j(k+1)+\tau c_j(k)=0$.
For $r\ge2$, put $h=k+r$.  The four terms in its finite column split as
\begin{align*}
 (v_r^{BC})^{\mathsf T}C_{\bullet,j}
 ={}&\bigl(c_j(h)+\tau^{2h-K}c_j(K-h)\bigr)\\
 &+\tau\bigl(c_j(h-1)+\tau^{2(h-1)-K}c_j(K-h+1)\bigr)=0.
\end{align*}
Thus the second kernel relation holds in both cases.

Because $K>2k$, the residues $j$ and $-j$ are disjoint and the first
$k$ rows give $C_I=I_k$.  By \eqref{eq:pairing-C} and $y=-q$,
\[
\begin{aligned}
 (D^C)^{\mathsf T}C&=A\diag((-1)^1,\ldots,(-1)^k),\\
 A_{ij}&=\sum_{y\in\Z}(-1)^{Ky}t^{Ky^2/2-jy}\Lambda_i(j-Ky).
\end{aligned}
\]
The sign $(-1)^{Ky}$ is $1$ for $C_0$ and $(-1)^y$ for $BC$.
Since $\det D^C_J=\zeta=(-1)^{nk}\omega$,
\eqref{eq:normalized-transfer} gives $P^\star_{n,k}=\det A$.
Finite row multilinearity proves the two identities.
\end{proof}

\subsection{The pair \texorpdfstring{$B_0$ and $B_-$}{B0 and B-}}

\begin{proposition}
Equations~\eqref{eq:main-B0} and~\eqref{eq:main-Bminus} hold for $n,k\ge1$.
\end{proposition}

\begin{proof}
Take $K=2k+1$ for $B_0$ and $K=2k$ for $B_-$, with $t=\tau^2$
in the latter case.  Define
\[
 b_j(m)=
 \begin{cases}
 (-1)^{Ky}t^{Ky^2/2-(j-1/2)y},&m=j-Ky,\quad y\in\Z,\\
 0,&m\not\equiv j\pmod K,
 \end{cases}
\]
and put
\[
 \gamma_j(m)=(-1)^m\bigl(b_j(m)+b_j(1-m)\bigr),
 \qquad C=\mathcal C^B(b).
\]
Changing $y$ to $y-1$ gives
$b_j(m+K)=(-1)^K t^{m+(K-1)/2}b_j(m)$.
Applying this relation to the two reflected terms yields
\begin{equation}\label{eq:ordinary-B-reflection}
 \gamma_j(1-m)=-\gamma_j(m),\qquad
 \gamma_j(m)=-t^{m-(K+1)/2}\gamma_j(K+1-m).
\end{equation}
For example, both $b_j(K+1-m)$ and $b_j(m-K)$ acquire the factor
$(-1)^K t^{(K+1)/2-m}$ from $b_j(1-m)$ and $b_j(m)$.
The parity in the definition of $\gamma_j$ leaves the minus sign
in \eqref{eq:ordinary-B-reflection} for either parity of $K$.

For $B_0$, the fixed index $m=k+1$ in the second reflection has
$\gamma_j(k+1)=0$; the remaining columns are killed by
$\gamma_j(k+r)+t^{r-1}\gamma_j(k+2-r)=0$.
For $B_-$, it gives instead
$\gamma_j(k+r)+\tau^{2r-1}\gamma_j(k+1-r)=0$.
The first reflection makes evaluation against $[m]_B$ agree with
$\gamma_j(m)$, so $V^{\mathsf T}C=0$ in both cases.

For $1\le p,j\le k$, $b_j(p)=\delta_{p,j}$, whereas
$b_j(1-p)=0$ because $1\le j+p-1\le2k-1<K$.
Thus $C_I=\diag((-1)^1,\ldots,(-1)^k)$ and $\det C_I=\omega$.
The common pairing \eqref{eq:pairing-B} gives
\[
 ((D^B)^{\mathsf T}C)_{ij}
 =A_{ij}:=\sum_{y\in\Z}
 (-1)^{Ky}t^{Ky^2/2-(j-1/2)y}h_i(j-Ky).
\]
Now \eqref{eq:normalized-transfer} and $\det D^B_J=\zeta$ give
$P^\star_{n,k}=\det A$.  For odd $K$ the affine sign is $(-1)^y$;
for even $K$ it is absent.  Finite multilinearity proves the stated
identities.
\end{proof}


\section{Fixed classes and adjacent residues}\label{sec:exceptional}

The remaining cases have one doubled complementary minor.  We treat
separately its three sources: a fixed residue, a fixed center, and the
overlap of adjacent residues.

\subsection{The fixed residue in \texorpdfstring{$C_t$}{Ct}}

\begin{proposition}
Equation~\eqref{eq:main-Ct} holds for $n,k\ge1$.
\end{proposition}

\begin{proof}
Set $K=2k$ and define
\[
 b_j(m)=
 \begin{cases}
 (-1)^q t^{kq^2+jq},&m=j+2kq,\quad q\in\Z,\\
 0,&m\not\equiv j\pmod{2k}.
 \end{cases}
\]
Let $c_j(m)=b_j(m)-b_j(-m)$ and $C=\mathcal C^C(b)$.
The quasiperiodicity is now
$b_j(m+2k)=-t^{m+k}b_j(m)$, so the calculation leading to
\eqref{eq:ordinary-C-reflection} gives
\[
 c_j(-m)=-c_j(m),\qquad
 c_j(m)=t^{m-k}c_j(2k-m).
\]
At $m=k+r$ this annihilates the column
$[k+r]_C-t^r[k-r]_C$; hence $(V^{C_t})^{\mathsf T}C=0$.

Unlike the ordinary cases, the residues $k$ and $-k$ coincide.
Their contributions add rather than cancel:
\[
 c_k(k+2kq)=2(-1)^qt^{kq(q+1)}\quad(q\in\Z).
\]
For $j<k$ the two residues remain disjoint.  Consequently,
$C_I=\diag(1,\ldots,1,2)$ and $\det C_I=2$.
This includes $k=1$, when the block consists of the single entry $2$.

Write
\[
\begin{gathered}
 a_{ij}(y)=t^{ky^2-jy}\Lambda_i(j-2ky),\\
 A^+_{ij}=\sum_y a_{ij}(y),\qquad
 A^-_{ij}=\sum_y(-1)^y a_{ij}(y).
\end{gathered}
\]
The pairing \eqref{eq:pairing-C}, with $y=-q$, gives
$(D^C)^{\mathsf T}C=A^-\diag((-1)^1,\ldots,(-1)^k)$.
Thus the fixed-residue minor and \eqref{eq:normalized-transfer} yield
$P^{C_t}_{n,k}=\tfrac12\det A^-$.

It remains to identify the even-sublattice expression, not the full
lattice sum.  The parity projector and finite multilinearity give
\[
 \sum_{y\in\latticeeven}\det[a_{ij}(y_i)]
 =\tfrac12\bigl(\det A^++\det A^-\bigr).
\]
In the last column, $y\mapsto1-y$ preserves $ky^2-ky$ and sends
$k-2ky$ to its negative.  Since $\Lambda_i$ is odd,
$a_{i,k}(1-y)=-a_{i,k}(y)$.  This fixed-point-free involution makes
the last column of $A^+$ zero.  The parity projection therefore equals
$\tfrac12\det A^-$, proving the identity.
\end{proof}

\subsection{The fixed center in \texorpdfstring{$B_+$}{B+}}

The sequence used here will also supply the adjacent-residue reflection
in the next subsection.  For $1\le j\le k$, set
\begin{equation}\label{eq:Bplus-sequence}
\begin{aligned}
 b_j(m)&=
 \begin{cases}
 (-1)^y t^{ky^2-(j-1)y},&m=j-2ky,\quad y\in\Z,\\
 0,&m\not\equiv j\pmod{2k},
 \end{cases}\\
 \Gamma_j(m)&=(-1)^m\bigl(b_j(m)+b_j(2-m)\bigr).
\end{aligned}
\end{equation}

\begin{lemma}\label{lem:center-reflection}
For every integer $m$,
\begin{align}
 b_j(m+2k)&=-t^{m+k-1}b_j(m),\notag\\
 \Gamma_j(2-m)&=\Gamma_j(m),\notag\\
 \Gamma_j(m)&=-t^{m-k-1}\Gamma_j(2k+2-m).
 \label{eq:center-reflection}
\end{align}
In particular, $\Gamma_j(k+1)=0$.
\end{lemma}

\begin{proof}
The first relation follows from $y\mapsto y-1$ in
\eqref{eq:Bplus-sequence}; the second is the defining symmetry.
For the last, the first relation gives
\[
 b_j(2k+2-m)=-t^{k+1-m}b_j(2-m),\qquad
 b_j(m-2k)=-t^{k+1-m}b_j(m).
\]
Since $m$ and $2k+2-m$ have the same parity, adding these relations
proves \eqref{eq:center-reflection}.  Its fixed index $m=k+1$
has eigenvalue $-1$ and hence vanishes.
\end{proof}

\Needspace{6\baselineskip}
\begin{proposition}
Equation~\eqref{eq:main-Bplus} holds for $n,k\ge1$.
\end{proposition}

\begin{proof}
Take $C=\mathcal C^+(b)$ for \eqref{eq:Bplus-sequence}.
The normalized center row is $C_{1,j}=-b_j(1)$, so
\[
 [1]_+^{\mathsf T}C_{\bullet,j}=2C_{1,j}
 =-2b_j(1)=\Gamma_j(1).
\]
Together with $\Gamma_j(2-m)=\Gamma_j(m)$, this shows that every
bracket $[m]_+$ evaluates to $\Gamma_j(m)$, including the fixed
center.  Lemma~\ref{lem:center-reflection} therefore gives
$(V^{B_+})^{\mathsf T}C=0$: the first column evaluates to
$\Gamma_j(k+1)$ and the others to
$\Gamma_j(k+r)+t^{r-1}\Gamma_j(k+2-r)=0$.

The first $k$ rows are
$C_I=\diag((-1)^1,\ldots,(-1)^k)$.
Indeed, $b_j(1)=\delta_{j,1}$; for $2\le p\le k$,
$b_j(p)=\delta_{p,j}$ and the reflected congruence would require
$2ky=j+p-2\in\{1,\ldots,2k-2\}$, which is impossible.
Thus $\det C_I=\omega$.  The factor $2$ occurs instead in
$\det D^+_J=2\zeta$, whose first diagonal coefficient is
$\dote_0+\dote_0$ by Lemma~\ref{lem:coefficient-kernels}.

By \eqref{eq:pairing-plus},
\[
 ((D^+)^{\mathsf T}C)_{ij}
 =A^{B_+}_{ij}:=\sum_y(-1)^yt^{ky^2-(j-1)y}h_i^+(j-2ky).
\]
In this pairing the center contributes exactly $h_i^+(1)b_j(1)$;
its coefficient is not doubled.  Formula~\eqref{eq:normalized-transfer}
now gives $P^{B_+}_{n,k}=\tfrac12\det A^{B_+}$, and finite row
multilinearity proves the result.
\end{proof}

\subsection{Adjacent residues in \texorpdfstring{$B_t$}{Bt}}

This case uses the same coefficient matrix $D^B$ as $B_0$ and $B_-$.
Its different row kernel satisfies
$\widehat h_i(m)=h_i(m)+h_i(m-1)$, directly from
$\ddote_r=\dote_r+\dote_{r-1}$.
Keep $b_j$ and $\Gamma_j$ from \eqref{eq:Bplus-sequence} and put
\[
\begin{aligned}
 a_j(m)&=b_j(m)+b_j(m+1),\\
 c_j(m)&=(-1)^m\bigl(a_j(m)+a_j(1-m)\bigr),\\
 \gamma_j(m)&=c_j(m)-tc_j(m-1).
\end{aligned}
\]
Thus $a_j$ is supported on the two adjacent residues $j$ and $j-1$,
and $c_j(1-m)=-c_j(m)$.

\begin{lemma}
The adjacent-residue sequences satisfy
\begin{align}
 c_j(m)&=\Gamma_j(m)-\Gamma_j(m+1),\label{eq:adjacent-difference}\\
 \gamma_j(m)&=(1+t)\Gamma_j(m)-\Gamma_j(m+1)-t\Gamma_j(m-1),
 \label{eq:adjacent-gamma}\\
 \gamma_j(m)&=-t^{m-k-1}\gamma_j(2k+2-m).
 \label{eq:adjacent-reflection}
\end{align}
In particular, $\gamma_j(k+1)=0$.
\end{lemma}

\begin{proof}
Expanding the definitions proves \eqref{eq:adjacent-difference}, and
subtracting $tc_j(m-1)$ gives \eqref{eq:adjacent-gamma}.
Let $T=2k+2$ and $d=m-k-1$.  Apply
\eqref{eq:center-reflection} at $m,m+1,m-1$ to the latter identity:
\begin{align*}
 \gamma_j(m)
 &=-t^d\bigl((1+t)\Gamma_j(T-m)
             -t\Gamma_j(T-m-1)-\Gamma_j(T-m+1)\bigr)\\
 &=-t^d\gamma_j(T-m).
\end{align*}
The fixed-index assertion follows at $m=k+1$.
\end{proof}

\begin{proposition}
Equation~\eqref{eq:main-Bt} holds for $n,k\ge1$.
\end{proposition}

\begin{proof}
Set $C=\mathcal C^B(a)$, so $C_{p,j}=c_j(p)$.
Evaluation against the first finite column gives
\[
 ([k+1]_B-t[k]_B)^{\mathsf T}C_{\bullet,j}
 =\gamma_j(k+1)=0.
\]
For $r\ge2$, the four terms in $u_r^B-t\widetilde u_r^B$ give
$\gamma_j(k+r)+t^{r-1}\gamma_j(k+2-r)=0$ by
\eqref{eq:adjacent-reflection}.  Hence
$(V^{B_t})^{\mathsf T}C=0$.

The first complementary block contains the overlap that fixes the
normalization.  For $1\le p,j\le k$, only $y=0$ contributes, and
\[
 a_j(p)=\delta_{p,j}+\delta_{p,j-1},\qquad
 a_j(1-p)=\delta_{p,1}\delta_{j,1}.
\]
The second equality follows because the two congruences require
$2ky=j+p-1$ or $2ky=j+p-2$; only the latter can occur, at $p=j=1$.
Therefore
\[
 C_{p,j}=(-1)^p\bigl(\delta_{p,j}+\delta_{p,j-1}
                         +\delta_{p,1}\delta_{j,1}\bigr)
 \quad(1\le p,j\le k).
\]
This block is upper bidiagonal, with diagonal $-2,1,-1,1,\ldots$,
and hence $\det C_I=2\omega$.  There is no doubling in
$\det D^B_J=\zeta$.

Finally, the common pairing \eqref{eq:pairing-B} and the definition of
$\widehat h_i$ give
\begin{align*}
 ((D^B)^{\mathsf T}C)_{ij}
 &=\sum_m h_i(m)\bigl(b_j(m)+b_j(m+1)\bigr)\\
 &=\sum_y(-1)^yt^{ky^2-(j-1)y}
          \bigl(h_i(j-2ky)+h_i(j-1-2ky)\bigr)\\
 &=\sum_y(-1)^yt^{ky^2-(j-1)y}\widehat h_i(j-2ky).
\end{align*}
The normalization formula yields one half of this determinant, and
finite multilinearity proves Theorem~\ref{thm:B}.
\end{proof}


\section{Width zero and classical specializations}\label{sec:boundary}

\subsection{The normalized empty rectangle}

\begin{proposition}
For every $n\ge1$, all seven normalized width-zero values are $1$.
Consequently, the five non-prefactored identities in
Theorem~\ref{thm:A} hold also at $k=0$.
\end{proposition}

\begin{proof}
The symmetric-group reduction \eqref{eq:signsum} is valid at $k=0$.
The alternant and row-folding calculation in Proposition~\ref{prop:finite}
therefore gives
\[
 \nu_\star(0)P^\star_{n,0}
 =\frac{\det(\Phi V^\star|_{k=0})}{\Delta},
\]
without using the complementary-kernel transfer.
Relative to the initial basis defining $\Delta$, these column
transformations are upper triangular.  Their diagonal coefficients in
column $r$ are
\begin{center}
\renewcommand{\arraystretch}{1.2}
\begin{tabular}{@{}lll@{}}
\toprule
Case & Diagonal coefficient & Product\\
\midrule
$C_0,BC,B_0$ & $1$ & $1$\\
$C_t,B_t$ & $1+t^r$ & $(-t;t)_n$\\
$B_-$ & $1-\tau^{2r-1}$ & $(\tau;\tau^2)_n$\\
$B_+$ & $1+t^{r-1}$ & $(-1;t)_n$\\
\bottomrule
\end{tabular}
\end{center}
For example, the $B_t$ columns are $(1+t)\ps_1$ and
\[
 (1+t^r)\ps_r-(t+t^{r-1})\ps_{r-1}\qquad(r\ge2),
\]
while the $BC$ columns are $\ph_1$ and
\[
 \ph_r+(\tau-\tau^{2r-1})\ph_{r-1}
        -\tau^{2r-2}\ph_{r-2}\qquad(r\ge2).
\]
For $B_+$ the first column is $C_0=2$ relative to $\chi_1=1$,
so its diagonal is indeed $1+t^0=2$.
Each product in the table equals $\nu_\star(0)$, proving the claim.
The five affine expressions without an external half-factor are $1$
by the empty-determinant convention.
\end{proof}

These are rational-function identities before cancellation; any
removable singularities of the displayed normalization factors are
understood after cancellation.  In particular, the conclusion does not
require a nonvanishing numerical value of every Pochhammer factor.

\subsection{The specialization \texorpdfstring{$t=0$}{t=0}}

The classical rectangular dual Jacobi--Trudi formulas are recovered in
three finite foldings.  They are the symplectic and orthogonal
determinants in \cite[Theorem~2.1]{AlbionEtAl2024}; see also the
lattice-path proofs in \cite{FulmekKrattenthaler1997}.

\begin{corollary}
For $n,k\ge1$, with $1\le a,r\le n$ on the left and
$1\le i,j\le k$ on the right,
\begin{align}
 \frac{\det[\ph_{k+r}(x_a)]}{\Delta_C(x)}
 &=\det[\dote_{n-i+j}-\dote_{n+i+j}],\label{eq:classical-C}\\
 \frac{\det[\ps_{k+r}(x_a)]}{\Delta_B(x)}
 &=\det[\dote_{n-i+j}+\dote_{n+i+j-1}]\notag\\
 &=\tfrac12\det[\ddote_{n-i+j}+\ddote_{n+i+j-1}],
 \label{eq:classical-B}\\
 \frac{\det[C_{k+r-1}(x_a)]}{\Delta_+(x)}
 &=\tfrac12\det[\dote_{n-i+j}+\dote_{n+i+j-2}].
 \label{eq:classical-plus}
\end{align}
The seven finite representations specialize to the corresponding left
sides: $C_0,C_t,BC$ to \eqref{eq:classical-C};
$B_0,B_-,B_t$ to \eqref{eq:classical-B}; and $B_+$ to
\eqref{eq:classical-plus}.
\end{corollary}

\Needspace{9\baselineskip}
\begin{proof}
At $t=0$ (and $\tau=0$ in the square-root cases), every finite
column in Table~\ref{tab:finite-columns} reduces to
$[k+r]_\diamond$.  This identifies the finite ratios.
In the affine formulas for $C_0,B_0,B_+,B_t$, every nonzero shift
has a strictly positive parameter exponent, leaving only $y=0$ in
each row and giving the displayed dual determinants.
For $C_t$ one should instead use its finite ratio: at $j=k,y=1$
the weight exponent is zero, so nonzero shifts do not disappear
entrywise.  The already proved parity identity identifies its
specialized even-sublattice sum with the same symplectic determinant.
\end{proof}


% GITHUB-EDITION:DISCLOSURE-BEGIN
\medskip
\noindent\textbf{Lean formalization.}
Theorems~A and~B, together with the auxiliary results listed in the
accompanying code, have been formalized in Lean~4. The source code,
theorem correspondence, and verification instructions are available at
\url{https://github.com/VerMoriarty/affine-dual-jacobi-trudi/tree/7d58ee5913d67d8036991a5a363260f24f0625bf}.

\Needspace{12\baselineskip}
\medskip
\noindent\textbf{AI use declaration.}
This Lean 4 formalization was developed by Hangzhou Logicon AI Studio, with AI assistance from OpenAI Codex (GPT-6.1 Sol), under human supervision and review.

Artificial intelligence was used to assist in exploring proof strategies and developing preliminary code for numerical experiments. All results were independently verified by the authors. The authors assume responsibility for all content.
\medskip
% GITHUB-EDITION:DISCLOSURE-END
\Needspace{30\baselineskip}
\begin{thebibliography}{99}
\raggedright

\bibitem{Aitken1931}
A.~C. Aitken,
\emph{Note on dual symmetric functions},
Proc. Edinburgh Math. Soc. (2) \textbf{2} (1931), no.~3, 164--167,
\texttt{doi:10.1017/S0013091500007719}.

\bibitem{AlbionEtAl2024}
S.~P. Albion, I.~Fischer, H.~H\"ongesberg and F.~Schreier-Aigner,
\emph{Skew symplectic and orthogonal characters through lattice paths},
European J. Combin. \textbf{122} (2024), Article 104000, 26 pages,
\texttt{doi:10.1016/j.ejc.2024.104000}.

\bibitem{FulmekKrattenthaler1997}
M.~Fulmek and C.~Krattenthaler,
\emph{Lattice path proofs for determinantal formulas for symplectic and
orthogonal characters},
J. Combin. Theory Ser. A \textbf{77} (1997), no.~1, 3--50,
\texttt{doi:10.1006/jcta.1996.2711}.

\bibitem{McDowell2023}
E.~McDowell,
\emph{Flagged Schur polynomial duality via a lattice path bijection},
Electron. J. Combin. \textbf{30} (2023), no.~1, Paper No.~P1.5, 17 pages,
\texttt{doi:10.37236/11200}.

\bibitem{MullerEtAl2016}
S.~M\"uller, E.~Feliu, G.~Regensburger, C.~Conradi, A.~Shiu and
A.~Dickenstein,
\emph{Sign conditions for injectivity of generalized polynomial maps with
applications to chemical reaction networks and real algebraic geometry},
Found. Comput. Math. \textbf{16} (2016), no.~1, 69--97,
\texttt{doi:10.1007/s10208-014-9239-3}.

\bibitem{Venkateswaran2013}
V.~Venkateswaran,
\emph{Symmetric and nonsymmetric Hall-Littlewood polynomials of type
$BC$},
arXiv:1209.2933v2 (2013).

\bibitem{Warnaar2026}
S.~Ole Warnaar,
\emph{Affine Jacobi--Trudi identities and $q,t$-Rogers--Ramanujan identities},
SIGMA \textbf{22} (2026), 062, 50 pages,
\texttt{doi:10.3842/SIGMA.2026.062}; arXiv:2511.17034v2.

\end{thebibliography}

\end{document}
